\documentclass[12pt,a4paper]{article}
\usepackage[utf8]{inputenc}
\usepackage{amsmath,amssymb,amsthm}
\usepackage[margin=1in]{geometry}

\newtheorem{theorem}{Theorem}
\newtheorem{lemma}[theorem]{Lemma}
\newtheorem{proposition}[theorem]{Proposition}
\newtheorem{corollary}[theorem]{Corollary}
\theoremstyle{definition}
\newtheorem{definition}[theorem]{Definition}

\title{Problem 20: Factorial Digit Sum}
\author{Project Euler Solutions}
\date{}

\begin{document}
\maketitle

\section{Problem Statement}
Let $S(m)$ denote the sum of the decimal digits of a positive integer~$m$. Compute~$S(100!)$.

\section{Formal Development}

\begin{theorem}[Digit count]
$d(n!) = \lfloor \log_{10}(n!) \rfloor + 1$, where
$\log_{10}(n!) = \sum_{k=1}^{n} \log_{10} k$.
\end{theorem}

\begin{proof}
$\log_{10}(n!) = \log_{10}\bigl(\prod_{k=1}^{n} k\bigr) = \sum_{k=1}^{n} \log_{10} k$.
The digit count of any positive integer~$m$ is $\lfloor \log_{10} m \rfloor + 1$.
\end{proof}

\begin{corollary}
$\log_{10}(100!) \approx 157.97$, so $100!$ has $158$ digits.
\end{corollary}

\begin{theorem}[Digit sum congruence]
$S(m) \equiv m \pmod{9}$ for every positive integer~$m$.
\end{theorem}

\begin{proof}
$10 \equiv 1 \pmod{9}$ implies $m = \sum a_k \cdot 10^k \equiv \sum a_k \pmod{9}$.
\end{proof}

\begin{theorem}[Legendre's formula]
For prime~$p$, $\nu_p(n!) = \sum_{i=1}^{\infty} \lfloor n/p^i \rfloor$.
\end{theorem}

\begin{proof}
Each integer in $\{1, \ldots, n\}$ contributes $\nu_p(k)$ to $\nu_p(n!)$.
Counting by powers: $\nu_p(n!) = \sum_{i \ge 1} \lfloor n/p^i \rfloor$.
\end{proof}

\begin{proposition}[$100! \equiv 0 \pmod{9}$]
$\nu_3(100!) = 33 + 11 + 3 + 1 = 48 \ge 2$, so $9 \mid 100!$ and $S(100!) \equiv 0 \pmod{9}$.
\end{proposition}

\begin{theorem}[Trailing zeros]
$\nu_5(n!) = \sum_{i \ge 1} \lfloor n/5^i \rfloor$ equals the number of trailing zeros of~$n!$.
\end{theorem}

\begin{proof}
Trailing zeros come from factors of $10 = 2 \times 5$. Since $\nu_2(n!) \ge \nu_5(n!)$
for all $n \ge 1$, the count equals $\nu_5(n!)$.
\end{proof}

\begin{corollary}
$\nu_5(100!) = 20 + 4 = 24$ trailing zeros, contributing nothing to $S(100!)$.
\end{corollary}

\paragraph{Verification.}
$S(100!) = 648 = 72 \times 9 \equiv 0 \pmod{9}$, consistent with $9 \mid 100!$.

\section{Complexity}

\begin{proposition}
$O(n^2 \log n)$ time, $\Theta(n \log n)$ space.
\end{proposition}

\begin{proof}
At step~$k$, the running product has $\Theta(k \log k)$ digits; multiplying by~$k$ costs
$O(k \log k)$. Summing: $\sum_{k=2}^{n} O(k \log k) = O(n^2 \log n)$.
Storage for the factorial requires $O(n \log n)$ digits.
\end{proof}

\section{Answer}

\[
\boxed{648}
\]

\end{document}
